\chapter{El núcleo y la radiactividad}\label{ch:g11:nucleus-radioactivity}

Ahora mismo, unos ocho mil \omterm{def:g11:fundamental-interactions:ladder}{núcleos} atómicos revientan dentro de tu
cuerpo cada segundo. Nada los ataca: ciertos \omterm{def:g11:fundamental-interactions:ladder}{núcleos} nacen inestables, y
cada uno se transforma tarde o temprano lanzando un fragmento a una
buena fracción de la velocidad de la luz. Este capítulo vuelve a la
planta de los $10^{-15}$ \omterm{def:g10:orders-of-magnitude:unit}{metros} del
\cref{ch:g11:fundamental-interactions} para ver qué \omterm{def:g11:fundamental-interactions:ladder}{núcleos} son
frágiles, cómo se rompen y cómo su impaciencia regular data cuevas y
vigila techos.

\section{Núcleos, nucleones e isótopos}

\begin{definition}[Notación nuclear]\label{def:g11:nucleus-radioactivity:notation}
Un \omterm{def:g11:fundamental-interactions:ladder}{núcleo} contiene $Z$ protones (de carga $+e$ cada uno) y $N$ neutrones
(sin carga), $A = Z + N$ nucleones en total; se escribe
${}^{A}_{Z}\mathrm{X}$, donde $\mathrm X$ es el símbolo químico, $Z$ el
\emph{número atómico}\index{número atómico} y $A$ el \emph{número
másico}\index{número másico}. Una especie determinada por $A$ y $Z$ es
un \emph{nucleido}\index{nucleido}; los nucleidos con el mismo $Z$ y
distinto $A$ son \emph{isótopos}\index{isótopo} de un mismo elemento:
la misma química (la nube de electrones solo ve $Ze$), distinto \omterm{def:g11:fundamental-interactions:ladder}{núcleo} y
distinta estabilidad.
\end{definition}

\begin{example}[Hidrógeno y carbono]\label{ex:g11:nucleus-radioactivity:hc}
El hidrógeno tiene tres \omterm{def:g11:nucleus-radioactivity:notation}{isótopos}: ${}^{1}_{1}\mathrm{H}$ (un protón
solo), ${}^{2}_{1}\mathrm{H}$ (deuterio: un protón y un neutrón) y el
inestable ${}^{3}_{1}\mathrm{H}$ (tritio: dos neutrones). El carbono
natural: en su mayoría ${}^{12}_{6}\mathrm{C}$ ($6$ protones, $6$
neutrones), un $1.1\,\%$ de ${}^{13}_{6}\mathrm{C}$ y un \omterm{def:g11:fundamental-interactions:ladder}{átomo} de cada
$10^{12}$ del inestable ${}^{14}_{6}\mathrm{C}$ ($8$ neutrones), héroe
del \cref{ex:g11:nucleus-radioactivity:dating}.
\end{example}

\section{Núcleos estables e inestables}

\begin{proposition}[El valle de estabilidad]\label{prop:g11:nucleus-radioactivity:valley}
Los \omterm{def:g11:nucleus-radioactivity:notation}{nucleidos} estables ocupan una banda estrecha del plano $(Z, N)$:
$N \approx Z$ para los \omterm{def:g11:fundamental-interactions:ladder}{núcleos} ligeros y después un exceso creciente de
neutrones, hasta $N \approx 1.5\,Z$. Todo \omterm{def:g11:nucleus-radioactivity:notation}{nucleido} fuera de la banda ---
con demasiados neutrones, con demasiados protones o demasiado grande
(más allá del plomo, $Z = 82$, ningún \omterm{def:g11:nucleus-radioactivity:notation}{nucleido} es verdaderamente
estable) --- se transforma tarde o temprano.
\end{proposition}
\admitted

\begin{remark}
Localizar la banda con exactitud es mecánica cuántica --- los volúmenes
universitarios. La pauta es el libro de cuentas del
\cref{ch:g11:fundamental-interactions}: el pegamento fuerte solo liga a
los vecinos y la repulsión de Coulomb abarca el \omterm{def:g11:fundamental-interactions:ladder}{núcleo} entero, así que
los \omterm{def:g11:fundamental-interactions:ladder}{núcleos} pesados necesitan neutrones de más como pegamento; pero un
exceso de neutrones también es inestable (de eso se encarga la
\omterm{def:g11:fundamental-interactions:weak}{interacción débil}).
\end{remark}

\begin{omfigure}
\begin{tikzpicture}[font=\small, x=0.048cm, y=0.033cm]
  \draw[-Stealth] (0,0) -- (108,0) node[below left=1pt] {$Z$};
  \draw[-Stealth] (0,0) -- (0,158) node[below right=1pt] {$N$};
  \draw (40,0) -- (40,-4) node[below] {$40$} (0,100) -- (-2,100) node[left] {$100$};
  \draw[dashed] (0,0) -- (105,105) node[right] {$N = Z$};
  \draw[omDef, line width=6pt, opacity=0.4] plot[smooth]
    coordinates {(1,1) (20,22) (40,50) (60,82) (83,126)};
  \draw[omDef, -Stealth] (52,30) node[below] {nucleidos estables} -- (57,72);
  \node at (40,120) {demasiados neutrones: $\beta^-$};
  \node at (82,18) {demasiados protones: $\beta^+$};
  \node[left] at (83,145) {demasiado pesado: $\alpha$};
\end{tikzpicture}

{\small Los \omterm{def:g11:nucleus-radioactivity:notation}{nucleidos} estables se pegan a $N = Z$ y después derivan
hacia el exceso de neutrones; por encima de la banda un \omterm{def:g11:fundamental-interactions:ladder}{núcleo} se
desintegra $\beta^-$, por debajo $\beta^+$ y, más allá del bismuto,
$\alpha$.}
\end{omfigure}

\begin{definition}[Radiactividad]\label{def:g11:nucleus-radioactivity:radioactivity}
La \emph{radiactividad}\index{radiactividad} es la transformación
espontánea de un \omterm{def:g11:fundamental-interactions:ladder}{núcleo} inestable en otro \omterm{def:g11:nucleus-radioactivity:notation}{nucleido}, con emisión de
radiación. Es \emph{aleatoria}: nada anuncia qué \omterm{def:g11:fundamental-interactions:ladder}{núcleo} se desintegrará
a continuación, y ninguna \omterm{def:g10:pressure:pressure}{presión}, temperatura ni química puede
acelerarla ni retrasarla; solo las poblaciones grandes son predecibles
(\cref{def:g11:nucleus-radioactivity:halflife}).
\end{definition}

\begin{remark}[Becquerel y los Curie]\label{rem:g11:nucleus-radioactivity:curie}
En 1896, Henri Becquerel descubrió que las sales de uranio velan una
placa fotográfica a través de papel negro, dentro de un cajón cerrado y
sin que se les aporte \omterm{def:g10:energy-conservation:energy}{energía}. Marie y Pierre Curie demostraron que el
efecto es atómico, lo midieron y aislaron el polonio y el radio, un
millón de veces más activos. La palabra \emph{\omterm{def:g11:nucleus-radioactivity:radioactivity}{radiactividad}} es de Marie
Curie; y también lo son dos premios Nobel.
\end{remark}

\section{Alfa, beta, gamma}

\begin{definition}[Las tres radiaciones históricas]\label{def:g11:nucleus-radioactivity:abg}
Tres radiaciones, bautizadas antes de que nadie supiera qué eran:
\begin{itemize}
  \item \emph{alfa}\index{radiación alfa} ($\alpha$): un \omterm{def:g11:fundamental-interactions:ladder}{núcleo} de helio
        ${}^{4}_{2}\mathrm{He}$ expulsado entero --- pesado, con carga
        doble, detenido por una hoja de papel o por unos centímetros de
        aire;
  \item \emph{beta}\index{radiación beta} ($\beta^-$): un electrón
        ${}^{\;0}_{-1}\mathrm{e}$ creado en el \omterm{def:g11:fundamental-interactions:ladder}{núcleo} y expulsado ---
        ligero, rápido, detenido por unos milímetros de aluminio;
  \item \emph{gamma}\index{radiación gamma} ($\gamma$): radiación
        electromagnética, mucho más energética que la luz --- nunca se
        detiene del todo, solo se
        \emph{atenúa}\index{atenuación}: un centímetro de plomo absorbe
        la mitad.
\end{itemize}
\end{definition}

\begin{omfigure}
\begin{tikzpicture}[font=\small]
  \fill[omExo] (-0.15,-0.3) rectangle (0,3.0); \node[left] at (-0.15,1.4) {fuente};
  \draw[omProp, thick] (2.2,-0.3) -- (2.2,3.0) node[above] {papel};
  \draw[omProp, line width=2.5pt] (4.4,-0.3) -- (4.4,3.0) node[above] {aluminio};
  \fill[omExo] (6.4,-0.3) rectangle (6.9,3.0); \node[above] at (6.65,3.0) {plomo};
  \draw[-Stealth, omDef, thick] (0,2.4) -- (2.15,2.4) node[midway, above] {$\alpha$};
  \draw[-Stealth, omDef, thick] (0,1.4) -- (4.35,1.4) node[pos=0.28, above] {$\beta^-$};
  \draw[omThm, thick, decorate,
        decoration={snake, amplitude=0.5mm, segment length=3mm}] (0,0.5) -- (6.4,0.5);
  \draw[-Stealth, omThm, thick, dashed] (6.9,0.5) -- (7.9,0.5);
  \node[omThm, above] at (1.1,0.62) {$\gamma$};
\end{tikzpicture}

{\small La $\alpha$ muere en el papel y la $\beta^-$ en unos milímetros
de aluminio; la $\gamma$ solo se adelgaza, incluso con plomo.}
\end{omfigure}

\begin{proposition}[Conservación en las ecuaciones nucleares]\label{prop:g11:nucleus-radioactivity:soddy}
En toda transformación nuclear, el \omterm{def:g11:nucleus-radioactivity:notation}{número másico} total $A$ y la carga
total --- la suma de los índices inferiores, contando $-1$ para un
electrón emitido --- son los mismos antes y después.
\end{proposition}
\admitted

\begin{remark}
La conservación de la carga ya gobernaba el
\cref{ch:g11:fundamental-interactions}; la conservación de $A$ dice que
los nucleones solo se convierten, nunca se crean ni se destruyen ---
contabilidad honrada para los volúmenes universitarios.
\end{remark}

\begin{method}[Ajustar una ecuación de desintegración]\label{met:g11:nucleus-radioactivity:balance}
\begin{enumerate}
  \item Escribir todos los actores con sus dos índices:
        $\alpha = {}^{4}_{2}\mathrm{He}$,
        $\beta^- = {}^{\;0}_{-1}\mathrm{e}$,
        $\beta^+ = {}^{0}_{+1}\mathrm{e}$; una $\gamma$ lleva
        $A = Z = 0$.
  \item Igualar las sumas de las $A$ y después las de las $Z$
        (\cref{prop:g11:nucleus-radioactivity:soddy}); resolver y leer el
        elemento hijo a partir de su $Z$ en la tabla \omterm{def:g10:signals-and-waves:period}{periódica}.
\end{enumerate}
\end{method}

\begin{example}[Una desintegración alfa]\label{ex:g11:nucleus-radioactivity:alpha}
El radio-226, el radio de los Curie, es un emisor $\alpha$:
${}^{226}_{88}\mathrm{Ra} \to {}^{A}_{Z}\mathrm{Y} + {}^{4}_{2}\mathrm{He}$
obliga a $A = 226 - 4 = 222$ y $Z = 88 - 2 = 86$; el elemento $86$ es el
radón, así que
${}^{226}_{88}\mathrm{Ra} \to {}^{222}_{86}\mathrm{Rn} + {}^{4}_{2}\mathrm{He}$
--- el gas radiactivo de los sótanos sin ventilar.
\end{example}

\begin{example}[Una desintegración beta menos]\label{ex:g11:nucleus-radioactivity:beta}
El carbono-14 está por encima de la banda ($8$ neutrones frente a $6$
protones):
${}^{14}_{6}\mathrm{C} \to {}^{14}_{7}\mathrm{N} + {}^{\;0}_{-1}\mathrm{e}$.
Comprobación: $14 = 14 + 0$ y $6 = 7 - 1$; un neutrón se ha convertido
en un protón más el electrón expulsado --- la \omterm{def:g11:fundamental-interactions:weak}{interacción débil} en
acción.
\end{example}

\begin{remark}[Beta más y gamma]\label{rem:g11:nucleus-radioactivity:betaplus}
Por debajo de la banda, la desintegración espejo $\beta^+$ convierte un
protón en un neutrón más un \emph{positrón}\index{positrón}
${}^{0}_{+1}\mathrm{e}$, la antipartícula del electrón:
${}^{18}_{9}\mathrm{F} \to {}^{18}_{8}\mathrm{O} + {}^{0}_{+1}\mathrm{e}$,
caballo de batalla de la tomografía por emisión de positrones. Y tras la
mayoría de las desintegraciones el \omterm{def:g11:fundamental-interactions:ladder}{núcleo} hijo, que nace tembloroso, se
asienta emitiendo además una $\gamma$.
\end{remark}

\begin{omfigure}
\begin{tikzpicture}[font=\small]
  \node (U)  at (0,0)   {${}^{238}_{92}\mathrm{U}$};
  \node (Th) at (2.6,0) {${}^{234}_{90}\mathrm{Th}$};
  \node (Pa) at (5.2,0) {${}^{234}_{91}\mathrm{Pa}$}; \node (U4) at (7.8,0) {${}^{234}_{92}\mathrm{U}$};
  \draw[-Stealth, omDef, thick] (U) -- (Th) node[midway, above] {$\alpha$};
  \draw[-Stealth, omProp, thick] (Th) -- (Pa) node[midway, above] {$\beta^-$};
  \draw[-Stealth, omProp, thick] (Pa) -- (U4) node[midway, above] {$\beta^-$};
\end{tikzpicture}

{\small Los primeros pasos de la cadena del uranio-238; once
desintegraciones más ($\alpha$ y $\beta^-$) terminan en el estable
${}^{206}_{82}\mathrm{Pb}$.}
\end{omfigure}

\section{Periodo de semidesintegración y actividad}

\begin{definition}[Periodo de semidesintegración]\label{def:g11:nucleus-radioactivity:halflife}
El \emph{periodo de semidesintegración}\index{periodo de
semidesintegración} $T_{1/2}$ de un \omterm{def:g11:nucleus-radioactivity:notation}{nucleido} es el tiempo tras el cual
se ha desintegrado la mitad de cualquier muestra grande. Cada \omterm{def:g10:signals-and-waves:period}{periodo}
adicional reduce a la mitad lo que queda: de $N_0$ \omterm{def:g11:fundamental-interactions:ladder}{núcleos} quedan
$N_0/2$ en $T_{1/2}$, $N_0/4$ en $2\,T_{1/2}$, $N_0/8$ en
$3\,T_{1/2}$ y $N_0/2^{\,n}$ tras $n$ \omterm{def:g10:signals-and-waves:period}{periodos}.
\end{definition}

\begin{remark}
Azar para un \omterm{def:g11:fundamental-interactions:ladder}{núcleo}, relojería para $10^{23}$: una muchedumbre inmensa
lanzando monedas, impredecible una a una, con la mitad eliminada en cada
ronda. Interpolar \emph{entre} \omterm{def:g10:signals-and-waves:period}{periodos} enteros exige la función
exponencial --- el curso que viene, con el volumen de matemáticas.
\end{remark}

\begin{example}[Veinte potencias de diez]\label{ex:g11:nucleus-radioactivity:range}
\begin{center}\small\begin{tabular}{llll}
polonio-214 & $\alpha$ & \qty{1.6e-4}{s} & eslabón de la cadena del radio \\
tecnecio-99m & $\gamma$ & $6.0$ horas & imagen médica \\
americio-241 & $\alpha$ & $432$ años & detectores de humo \\
carbono-14 & $\beta^-$ & \num{5700} años & arqueología \\
potasio-40 & $\beta$ & \num{1.25e9} años & en todos los plátanos \\
uranio-238 & $\alpha$ & \num{4.5e9} años & calor interno de la Tierra \\
\end{tabular}\end{center}
\end{example}

\begin{example}[Datación por carbono-14]\label{ex:g11:nucleus-radioactivity:dating}
La materia viva intercambia carbono con la atmósfera, así que mantiene
la proporción atmosférica de carbono-14; al morir cesa la entrada y la
proporción se reduce a la mitad cada \num{5700} años. Un carbón de una
cueva pintada muestra la cuarta parte de la proporción de los vivos: un
cuarto es la mitad de la mitad, luego el fuego ardió hace dos \omterm{def:g10:signals-and-waves:period}{periodos},
$2 \times \num{5700} = \num{11400}$ años.
\end{example}

\begin{definition}[Actividad]\label{def:g11:nucleus-radioactivity:activity}
La \emph{actividad}\index{actividad} $\mathcal A$ de una muestra es su
número de desintegraciones por segundo, medida en
\emph{becquerelios}\index{becquerelio}: $\qty{1}{Bq} = 1$
desintegración por segundo. Cuantos menos \omterm{def:g11:fundamental-interactions:ladder}{núcleos} sobreviven, menos
desintegraciones hay, así que la actividad también se reduce a la mitad
en cada \omterm{def:g10:signals-and-waves:period}{periodo}.
\end{definition}

\begin{omfigure}
\begin{tikzpicture}[font=\small, x=1.1cm, y=3.0cm]
  \draw[-Stealth] (0,0) -- (5.7,0) node[below left=1pt] {$t$ (en $T_{1/2}$)};
  \draw[-Stealth] (0,0) -- (0,1.15) node[right] {$\mathcal A/\mathcal A_0$};
  \foreach \n/\v in {0/1, 1/0.5, 2/0.25, 3/0.125, 4/0.0625}{
    \filldraw[fill=omDef!35, draw=omDef] (\n+0.15,0) rectangle (\n+0.85,\v);
    \node[below] at (\n+0.5,0) {$\n$};}
  \node[left] at (0,1) {$1$}; \node[left] at (0,0.5) {$\tfrac12$};
  \node[left] at (0,0.25) {$\tfrac14$};
  \draw[dashed, omExo] (0,0.5) -- (1.15,0.5) (0,0.25) -- (2.15,0.25);
\end{tikzpicture}

{\small La escalera de las mitades: tras $n$ \omterm{def:g10:signals-and-waves:period}{periodos} la \omterm{def:g11:nucleus-radioactivity:activity}{actividad} vale
$\mathcal A_0/2^{\,n}$ --- una octava parte tras tres, una milésima tras
diez.}
\end{omfigure}

\begin{example}[Órdenes de magnitud]\label{ex:g11:nucleus-radioactivity:orders}
Todo es un poco radiactivo. Un plátano: unos \qty{15}{Bq}
(potasio-40). Un cuerpo humano: unos \qty{8000}{Bq} --- la entradilla de
este capítulo, unos quinientos plátanos. Un \omterm{def:g10:orders-of-magnitude:unit}{metro} cúbico de granito:
alrededor de \qty{3e6}{Bq} (uranio y su cadena, de donde sale el radón
de los sótanos). Una inyección de tecnecio: unos \qty{5e8}{Bq}, que
desaparecen en pocos días.
\end{example}

\section{Servidores y peligros}

\begin{example}[Tres carreras de un núcleo inestable]\label{ex:g11:nucleus-radioactivity:uses}
\emph{Medicina}: el tecnecio-99m, emisor $\gamma$ puro con
$T_{1/2} = 6$ horas, se fija a una \omterm{def:g11:fundamental-interactions:ladder}{molécula} que absorbe el órgano
buscado; una cámara filma desde fuera las $\gamma$ y, al día siguiente,
el trazador casi ha desaparecido --- mientras que la radioterapia
invierte la lógica y concentra haces para destruir un tumor.
\emph{Datación}: carbono-14 para madera, hueso y tejido hasta unos
\num{50000} años; uranio-238 para las rocas --- y para la Tierra, con
sus \num{4.5e9} años. \emph{Detectores de humo}: una mota de
americio-241 ioniza el aire de una pequeña cámara y deja pasar una
corriente diminuta; el humo ahoga la corriente y salta la alarma --- y
las $\alpha$ mueren en unos centímetros de aire.
\end{example}

\begin{remark}[La dosis y el sievert]\label{rem:g11:nucleus-radioactivity:dose}
La radiación arranca electrones de las \omterm{def:g11:fundamental-interactions:ladder}{moléculas} y puede dañar el ADN.
El daño biológico se sigue con la \emph{dosis}\index{dosis}, en
\emph{sieverts}\index{sievert} (\unit{Sv}), que pondera la \omterm{def:g10:energy-conservation:energy}{energía}
depositada por \omterm{def:g10:orders-of-magnitude:unit}{kilogramo} de tejido según lo dañina que sea cada
radiación. Tres reglas: distancia, blindaje (para la $\alpha$ basta la
piel --- salvo que el emisor se \emph{inhale}, que es el crimen del
radón; para la $\beta$, aluminio; para la $\gamma$, plomo) y tiempo. La
dosimetría propiamente dicha es materia universitaria.
\end{remark}

\section{Ejercicios}

\begin{exercise}[$\star$]\label{exo:g11:nucleus-radioactivity:1}
Da el número de protones y de neutrones de ${}^{16}_{8}\mathrm{O}$,
${}^{27}_{13}\mathrm{Al}$, ${}^{60}_{27}\mathrm{Co}$,
${}^{235}_{92}\mathrm{U}$ y ${}^{3}_{1}\mathrm{H}$.
\end{exercise}

\begin{exercise}[$\star$]\label{exo:g11:nucleus-radioactivity:2}
Entre ${}^{14}_{6}\mathrm{C}$, ${}^{14}_{7}\mathrm{N}$,
${}^{12}_{6}\mathrm{C}$, ${}^{13}_{6}\mathrm{C}$ y
${}^{16}_{8}\mathrm{O}$, ¿cuáles son \omterm{def:g11:nucleus-radioactivity:notation}{isótopos} entre sí? ¿Por qué
${}^{14}_{6}\mathrm{C}$ y ${}^{14}_{7}\mathrm{N}$ \emph{no} son
\omterm{def:g11:nucleus-radioactivity:notation}{isótopos}, pese a tener el mismo \omterm{def:g11:nucleus-radioactivity:notation}{número másico}?
\end{exercise}

\begin{exercise}[$\star$]\label{exo:g11:nucleus-radioactivity:3}
El polonio-210, ${}^{210}_{84}\mathrm{Po}$, es un emisor $\alpha$.
Escribe la ecuación de desintegración y nombra al hijo ($Z = 81$: talio;
$82$: plomo; $83$: bismuto).
\end{exercise}

\begin{exercise}[$\star$]\label{exo:g11:nucleus-radioactivity:4}
Escribe las ecuaciones de desintegración $\beta^-$ del cobalto-60
(${}^{60}_{27}\mathrm{Co}$) y del yodo-131
(${}^{131}_{53}\mathrm{I}$); hijos: $Z = 28$, níquel; $Z = 54$, xenón.
\end{exercise}

\begin{exercise}[$\star$]\label{exo:g11:nucleus-radioactivity:5}
Una fuente de yodo-131 ($T_{1/2} = 8.0$ días) tiene hoy una \omterm{def:g11:nucleus-radioactivity:activity}{actividad} de
\qty{800}{MBq}. ¿Qué \omterm{def:g11:nucleus-radioactivity:activity}{actividad} tiene al cabo de $16$ días? ¿Y de $32$
días?
\end{exercise}

\begin{exercise}[$\star\star$]\label{exo:g11:nucleus-radioactivity:6}
Cuatro radiaciones atraviesan un \omterm{def:g11:magnetic-fields:bfield}{campo magnético} intenso y después unas
pantallas. Identifica cada una: (a) la detiene el papel y el \omterm{def:g11:electric-gravitational-fields:field}{campo}
apenas la desvía; (b) atraviesa el papel, la detienen \qty{4}{mm} de
aluminio y el \omterm{def:g11:electric-gravitational-fields:field}{campo} la desvía mucho; (c) atraviesa los dos, el plomo
solo la \omterm{def:g11:nucleus-radioactivity:abg}{atenúa} y no se desvía; (d) como (b), pero desviada hacia el otro
lado.
\end{exercise}

\begin{exercise}[$\star\star$]\label{exo:g11:nucleus-radioactivity:7}
El radón-222 (${}^{222}_{86}\mathrm{Rn}$) abre una cadena rápida: una
desintegración $\alpha$, después otra $\alpha$ y después una $\beta^-$.
Escribe las tres ecuaciones ($Z = 82$: plomo; $83$: bismuto; $84$:
polonio).
\end{exercise}

\begin{exercise}[$\star\star$]\label{exo:g11:nucleus-radioactivity:8}
El flúor-18, ${}^{18}_{9}\mathrm{F}$, es el trazador de la tomografía
por emisión de positrones; el flúor estable es
${}^{19}_{9}\mathrm{F}$. ¿En qué lado de la banda de estabilidad está?
Predice su modo de desintegración y escribe la ecuación ($Z = 8$:
oxígeno).
\end{exercise}

\begin{exercise}[$\star\star$]\label{exo:g11:nucleus-radioactivity:9}
Un paciente recibe \qty{500}{MBq} de tecnecio-99m
($T_{1/2} = 6.0$ horas) a las 08:00. ¿Qué \omterm{def:g11:nucleus-radioactivity:activity}{actividad} hay a las 08:00 de
la mañana siguiente? ¿Al cabo de cuántas horas baja por primera vez de
\qty{1}{MBq}?
\end{exercise}

\begin{exercise}[$\star\star$]\label{exo:g11:nucleus-radioactivity:10}
Un hueso de una turbera muestra una proporción de carbono-14 igual a la
octava parte de la de un hueso vivo ($T_{1/2} = \num{5700}$ años).
¿Cuántos años tiene el hueso? ¿Y por qué el carbono-14 no sirve para
datar huesos de dinosaurio (de unos $10^{8}$ años)?
\end{exercise}

\begin{exercise}[$\star\star$]\label{exo:g11:nucleus-radioactivity:11}
Con el valle de estabilidad, predice el modo de desintegración
($\alpha$, $\beta^-$ o $\beta^+$) de ${}^{14}_{6}\mathrm{C}$,
${}^{11}_{6}\mathrm{C}$ y ${}^{238}_{92}\mathrm{U}$, justificando cada
uno en una línea (carbono estable: ${}^{12}_{6}\mathrm{C}$ y
${}^{13}_{6}\mathrm{C}$).
\end{exercise}

\begin{exercise}[$\star\star\star$]\label{exo:g11:nucleus-radioactivity:12}
La cadena del uranio-238 termina, muchos pasos después, en el estable
${}^{206}_{82}\mathrm{Pb}$; cada paso es una $\alpha$ o una $\beta^-$.
Usando solo la conservación de $A$, halla el número de pasos $\alpha$; y
después, con la conservación de $Z$, el número de pasos $\beta^-$.
\end{exercise}

\begin{exercise}[$\star\star\star$]\label{exo:g11:nucleus-radioactivity:13}
La fuente de cobalto-60 de un hospital ($T_{1/2} = 5.3$ años) tiene una
\omterm{def:g11:nucleus-radioactivity:activity}{actividad} de \qty{6.4}{TBq}; la normativa permite deshacerse de ella por
debajo de \qty{0.1}{TBq}. ¿Al cabo de cuántos años se puede retirar?
\end{exercise}

\begin{exercise}[$\star\star\star$]\label{exo:g11:nucleus-radioactivity:14}
La \omterm{def:g11:nucleus-radioactivity:activity}{actividad} de tu cuerpo es de unos \qty{8000}{Bq}. ¿Cuántos \omterm{def:g11:fundamental-interactions:ladder}{núcleos}
tuyos se desintegran al día? ¿Y a lo largo de una vida de $80$ años? Un
plátano añade unos \qty{15}{Bq} mientras lo digieres: comenta en una
frase los titulares que temen a cada \omterm{def:g11:nucleus-radioactivity:activity}{becquerelio}.
\end{exercise}

\begin{exercise}[$\star\star\star$]\label{exo:g11:nucleus-radioactivity:15}
Un detector de humo contiene \qty{0.2}{\micro\gram} de americio-241
($T_{1/2} = 432$ años), con una \omterm{def:g11:nucleus-radioactivity:activity}{actividad} de unos \qty{25}{kBq}.
¿Cuántas desintegraciones son al día? Tras diez años de servicio, ¿el
americio se ha desintegrado mucho menos de la mitad, aproximadamente la
mitad o mucho más --- y es el agotamiento de la fuente la razón de que
se cambien los detectores? Y por último: ¿por qué esta fuente $\alpha$ es
inofensiva en el techo y peligrosa si se tritura la cápsula y se inhala
el polvo?
\end{exercise}

\section{Problema: la momia, el reactor y el detector de humo}

\begin{problem}[{Problema de fin de semana --- tres carreras de un
núcleo inestable: una momia datada a base de mitades, el árbol
genealógico del uranio ajustado, un detector de humo auditado y por qué
las estrellas tenían que ser mortales}]
\label{pb:g11:nucleus-radioactivity:1}
Tres \omterm{def:g11:fundamental-interactions:ladder}{núcleos} inestables, tres oficios: el carbono-14 pone fecha a toda
planta y todo animal muerto, el uranio-238 calienta un planeta desde
dentro y el americio-241 vigila techos. Datos:
$T_{1/2} = \num{5700}$ años (C-14), \num{4.5e9} años (U-238) y $432$
años (Am-241); el carbono vivo muestra $13.6$ desintegraciones por
minuto y gramo; \omterm{def:g11:nucleus-radioactivity:notation}{números atómicos}: N $7$, O $8$, Pb $82$, Bi $83$, Po
$84$, Rn $86$, Ra $88$, Th $90$, Pa $91$, U $92$, Np $93$, Am $95$.

\medskip
\noindent\textbf{Parte I --- La momia.}
\begin{enumerate}
  \item Da la composición (protones y neutrones) de
        ${}^{12}_{6}\mathrm{C}$ y ${}^{14}_{6}\mathrm{C}$. ¿Por qué
        tienen una química idéntica dentro de un cuerpo?
  \item ¿Por qué es constante la proporción de carbono-14 en un cuerpo
        vivo, y por qué empieza a bajar al morir?
  \item Escribe la ecuación de desintegración del carbono-14. ¿Escaparía
        su radiación a través de las vendas de lino de la momia?
  \item Un gramo de carbono del lino muestra $6.8$ desintegraciones por
        minuto. ¿Qué fracción del ritmo de los vivos es? Data la momia.
  \item Una segunda ``momia'', joya de una colección privada, muestra
        $13.2$ desintegraciones por minuto y gramo. ¿Veredicto?
  \item Estima el ritmo de una muestra de diez \omterm{def:g10:signals-and-waves:period}{periodos} de antigüedad y
        explica por qué la datación por carbono se apaga más allá de unos
        \num{57000} años.
\end{enumerate}

\medskip
\noindent\textbf{Parte II --- El reactor bajo el prado.}
\begin{enumerate}[resume]
  \item Da la composición de ${}^{238}_{92}\mathrm{U}$. ¿Por qué
        necesitan los \omterm{def:g11:fundamental-interactions:ladder}{núcleos} pesados semejante excedente de neutrones?
  \item Escribe su ecuación de desintegración $\alpha$.
  \item El hijo se desintegra después $\beta^-$, y el nieto otra vez
        $\beta^-$. Escribe las dos ecuaciones. ¿Qué elemento reaparece?
  \item La cadena termina en ${}^{206}_{82}\mathrm{Pb}$. Cuenta sus
        pasos $\alpha$ y después sus pasos $\beta^-$.
  \item La Tierra se formó hace \num{4.5e9} años. ¿Qué fracción de su
        uranio-238 primordial queda hoy? ¿Y dentro de $5$ mil millones
        de años más?
  \item En una frase: ¿qué alimenta en la superficie ese calor de
        desintegración enterrado? (Piensa en los volcanes y en la deriva
        de los continentes.)
\end{enumerate}

\medskip
\noindent\textbf{Parte III --- El detector de humo.}
\begin{enumerate}[resume]
  \item Escribe la ecuación de desintegración $\alpha$ del
        ${}^{241}_{95}\mathrm{Am}$.
  \item Su \omterm{def:g11:nucleus-radioactivity:activity}{actividad} es de \qty{33}{kBq}: ¿cuántas desintegraciones por
        segundo? ¿Y al día?
  \item Explica el detector: ¿qué les hacen las partículas $\alpha$ al
        aire de la cámara, y qué cambia el humo?
  \item ¿Por qué es un emisor $\alpha$ la elección correcta aquí --- y
        un emisor $\gamma$ de la misma \omterm{def:g11:nucleus-radioactivity:activity}{actividad} la peor de todas?
  \item Tras un \omterm{def:g11:nucleus-radioactivity:halflife}{periodo de semidesintegración} (¡$432$ años!) la
        \omterm{def:g11:nucleus-radioactivity:activity}{actividad} sería de \qty{16.5}{kBq}; el manual pide sustituir el
        detector al cabo de diez años. ¿Es la fuente el eslabón débil?
\end{enumerate}

\medskip
\noindent\textbf{Parte IV --- Por qué las estrellas tenían que ser
mortales.}
\begin{enumerate}[resume]
  \item Las estrellas moribundas forjaron y esparcieron todos los
        elementos más pesados que el helio. ¿Qué fracción del uranio-238
        forjado hace \num{9e9} años sobrevive hoy? ¿Por qué la abundancia
        que le queda a la Tierra hace del planeta algo mucho más joven
        que las estrellas más viejas?
  \item El calor interno radiactivo de la Tierra impulsa el vulcanismo y
        la tectónica de placas, y mantiene el \omterm{def:g11:fundamental-interactions:ladder}{núcleo} agitado fabricando
        el escudo magnético del \cref{ch:g11:magnetic-fields}. ¿Cómo
        sería un planeta hecho solo de \omterm{def:g11:fundamental-interactions:ladder}{átomos} estables?
  \item Una línea, como poesía y como contabilidad: ¿dónde se forjaron
        tus \omterm{def:g11:fundamental-interactions:ladder}{átomos} y qué sigue calentando el suelo que pisas?
\end{enumerate}
\end{problem}
